
\Data\ID{1003124202}
\Updated{2020-07-15 03:07:21}
\Level{A}
\SubArea{Absolute value}
\LANGen{
\Question{
The value of \( |2-6|-|-1+3| \)  is equal to:}
{\( 2 \)}{\( -6 \)}{\( 0 \)}{\( 8 \)}{}{}}
\LANGpl{
\Question{
Liczba \( |2-6|-|-1+3| \)  jest równa:}
{\( 2 \)}{\( -6 \)}{\( 0 \)}{\( 8 \)}{}{}}
\LANGcs{
\Question{
Hodnota \( |2-6|-|-1+3| \)  se rovná:}
{\( 2 \)}{\( -6 \)}{\( 0 \)}{\( 8 \)}{}{}}
\LANGsk{
\Question{
Hodnota \( |2-6|-|-1+3| \)  je rovná:}
{\( 2 \)}{\( -6 \)}{\( 0 \)}{\( 8 \)}{}{}}
\LANGes{
\Question{
El valor de \( |2-6|-|-1+3| \)  equivale a:}
{\( 2 \)}{\( -6 \)}{\( 0 \)}{\( 8 \)}{}{}}
\End

\Data\ID{1003124206}
\Updated{2021-01-05 08:01:04}
\Level{A}
\SubArea{Absolute value}
\LANGen{
\Question{
Simplifying \( \left|\sqrt5-3\right|-\left|2\sqrt5-4\right| \) we get:}
{\( -3\sqrt5+7 \)}{\( -\sqrt5+1 \)}{\( -3\sqrt5+1 \)}{\( -\sqrt5+7 \)}{}{}}
\LANGpl{
\Question{
Liczba  \( \left|\sqrt5-3\right|-\left|2\sqrt5-4\right| \) jest równa:}
{\( -3\sqrt5+7 \)}{\( -\sqrt5+1 \)}{\( -3\sqrt5+1 \)}{\( -\sqrt5+7 \)}{}{}}
\LANGcs{
\Question{
Zjednodušením \( \left|\sqrt5-3\right|-\left|2\sqrt5-4\right| \) dostaneme:}
{\( -3\sqrt5+7 \)}{\( -\sqrt5+1 \)}{\( -3\sqrt5+1 \)}{\( -\sqrt5+7 \)}{}{}}
\LANGsk{
\Question{
Zjednodušením \( \left|\sqrt5-3\right|-\left|2\sqrt5-4\right| \) dostaneme:}
{\( -3\sqrt5+7 \)}{\( -\sqrt5+1 \)}{\( -3\sqrt5+1 \)}{\( -\sqrt5+7 \)}{}{}}
\LANGes{
\Question{
Simplificando \( \left|\sqrt5-3\right|-\left|2\sqrt5-4\right| \) obtenemos:}
{\( -3\sqrt5+7 \)}{\( -\sqrt5+1 \)}{\( -3\sqrt5+1 \)}{\( -\sqrt5+7 \)}{}{}}
\End

\Data\ID{1003187002}
\Updated{2020-07-15 03:07:21}
\Level{A}
\SubArea{Absolute value}
\LANGen{
\Question{
Evaluate the expression \( \left|\left(1-\sqrt2\right)^2 \right|+\left|\left(1+\sqrt2\right)^2\right|-|-6| \).}
{\( 0 \)}{\( 12 \)}{\( 4\sqrt2 \)}{\( -4 \)}{}{}}
\LANGpl{
\Question{
Wartością wyrażenia \( \left|\left(1-\sqrt2\right)^2 \right|+\left|\left(1+\sqrt2\right)^2\right|-|-6| \) jest:}
{\( 0 \)}{\( 12 \)}{\( 4\sqrt2 \)}{\( -4 \)}{}{}}
\LANGcs{
\Question{
Vypočti hodnotu výrazu \( \left|\left(1-\sqrt2\right)^2 \right|+\left|\left(1+\sqrt2\right)^2\right|-|-6| \).}
{\( 0 \)}{\( 12 \)}{\( 4\sqrt2 \)}{\( -4 \)}{}{}}
\LANGsk{
\Question{
Vypočítaj hodnotu výrazu \( \left|\left(1-\sqrt2\right)^2 \right|+\left|\left(1+\sqrt2\right)^2\right|-|-6| \).}
{\( 0 \)}{\( 12 \)}{\( 4\sqrt2 \)}{\( -4 \)}{}{}}
\LANGes{
\Question{
Calcula la expresión \( \left|\left(1-\sqrt2\right)^2 \right|+\left|\left(1+\sqrt2\right)^2\right|-|-6| \).}
{\( 0 \)}{\( 12 \)}{\( 4\sqrt2 \)}{\( -4 \)}{}{}}
\End

\Data\ID{9000078509}
\Updated{2021-01-05 07:01:30}
\Level{A}
\SubArea{Absolute value}
\LANGen{
\Question{
Evaluate the following expression. 
\[ 
 |3 - 7|-|2(-4)| + |(-5)(-2)|
\]}
{\(6\)}{\(14\)}{\(22\)}{\(- 2\)}{}{}}
\LANGpl{
\Question{
Oblicz wyrażenie.
\[ 
 |3 - 7|-|2(-4)| + |(-5)(-2)|
\]}
{\(6\)}{\(14\)}{\(22\)}{\(- 2\)}{}{}}
\LANGcs{
\Question{
Určete hodnotu výrazu \(|3 - 7|-|2(-4)| + |(-5)(-2)|\).}
{\(6\)}{\(14\)}{\(22\)}{\(- 2\)}{}{}}
\LANGsk{
\Question{
Určte hodnotu výrazu \(|3 - 7|-|2(-4)| + |(-5)(-2)|\).}
{\(6\)}{\(14\)}{\(22\)}{\(- 2\)}{}{}}
\LANGes{
\Question{
Evalúa la siguiente expresión.
\[ 
 |3 - 7|-|2(-4)| + |(-5)(-2)|
\]}
{\(6\)}{\(14\)}{\(22\)}{\(- 2\)}{}{}}
\End

\Data\ID{1003124201}
\Updated{2021-01-05 08:01:44}
\Level{B}
\SubArea{Absolute value}
\LANGen{
\Question{
Which equation describes real numbers \( x \) that are equidistant from the numbers \( 6 \) and \( -3 \) on the number line?}
{\( |x-6|=|x+3| \)}{\( |x+6|=|x+3| \)}{\( |x-6|=|x-3| \)}{\( |x+6|=|x-3| \)}{}{}}
\LANGpl{
\Question{
Zbiór liczb \( x \), które na osi liczbowej są równo odległe od liczb  \( 6 \) i \( -3 \) można opisać za pomocą równania:}
{\( |x-6|=|x+3| \)}{\( |x+6|=|x+3| \)}{\( |x-6|=|x-3| \)}{\( |x+6|=|x-3| \)}{}{}}
\LANGcs{
\Question{
Reálná čísla \( x \), jejichž vzdálenost od čísel \( 6 \) a \( -3 \) na číselné ose je stejná, vyhovují rovnici:}
{\( |x-6|=|x+3| \)}{\( |x+6|=|x+3| \)}{\( |x-6|=|x-3| \)}{\( |x+6|=|x-3| \)}{}{}}
\LANGsk{
\Question{
Reálne čísla \( x\), ktoré sú na číselnej osi rovnako vzdialené od čísel \( 6 \) a \( -3 \) môžme opísať rovnicou?}
{\( |x-6|=|x+3| \)}{\( |x+6|=|x+3| \)}{\( |x-6|=|x-3| \)}{\( |x+6|=|x-3| \)}{}{}}
\LANGes{
\Question{
Los números reales \( x \), cuya distancia a los números \( 6 \) y \( -3 \) es la misma en la recta numérica, equivalen a la ecuación:}
{\( |x-6|=|x+3| \)}{\( |x+6|=|x+3| \)}{\( |x-6|=|x-3| \)}{\( |x+6|=|x-3| \)}{}{}}
\End

\Data\ID{1003124203}
\Updated{2020-11-30 04:11:24}
\Level{B}
\SubArea{Absolute value}
\LANGen{
\Question{
Assuming \( x < 0 \), the expression  \( \bigl| |x|+2 \bigr| \)  is equal to:}
{\( -x+2 \)}{\( x+2  \)}{\( -x-2  \)}{\( x-2 \)}{}{}}
\LANGpl{
\Question{
Dla \( x < 0 \), wyrażenie  \( \bigl| |x|+2 \bigr| \)  jest równe:}
{\( -x+2 \)}{\( x+2  \)}{\( -x-2  \)}{\( x-2 \)}{}{}}
\LANGcs{
\Question{
Předpokládejme, že \( x < 0 \). Výraz \( \bigl| |x|+2 \bigr| \) se rovná:}
{\( -x+2 \)}{\( x+2  \)}{\( -x-2  \)}{\( x-2 \)}{}{}}
\LANGsk{
\Question{
Predpokladajme, že \( x < 0 \). Výraz  \( \bigl| |x|+2 \bigr| \)  sa rovná:}
{\( -x+2 \)}{\( x+2  \)}{\( -x-2  \)}{\( x-2 \)}{}{}}
\LANGes{
\Question{
Suponiendo que \( x < 0 \), la expresión  \( \bigl| |x|+2 \bigr| \)  equivale a:}
{\( -x+2 \)}{\( x+2  \)}{\( -x-2  \)}{\( x-2 \)}{}{}}
\End

\Data\ID{1003124204}
\Updated{2020-07-15 03:07:21}
\Level{B}
\SubArea{Absolute value}
\LANGen{
\Question{
Let \( x\neq0 \). Complete the following sentence to get a true statement. The solution set of the inequality \( \frac{|x|}x>2 \)}
{does not contain any integer.}{contains \( 2 \) integers.}{contains only natural numbers.}{contains infinitely many integers.}{}{}}
\LANGpl{
\Question{
Wiadomo, że  \( x\neq0 \). Zatem do zbioru rozwiązań nierówności \( \frac{|x|}x>2 \)}
{nie należy żadna liczba całkowita.}{należą \( 2 \) liczby całkowite.}{należą tylko liczby naturalne.}{należy nieskończenie wiele liczb całkowitych.}{}{}}
\LANGcs{
\Question{
Nechť \( x\neq0 \). Doplň následující větu tak, aby tvrzení bylo pravdivé. Množina řešení nerovnice \( \frac{|x|}x>2 \)}
{neobsahuje žádné celé číslo.}{obsahuje \( 2 \) celá čísla.}{obsahuje jen přirozená čísla.}{obsahuje nekonečně mnoho celých čísel.}{}{}}
\LANGsk{
\Question{
Nech \( x\neq0 \). Doplň následujúcu vetu tak, aby tvrdenie bolo pravdivé. Množina riešenie nerovnice \( \frac{|x|}x>2 \)}
{neobsahuje žiadne celé číslo.}{obsahuje \( 2 \) celé čísla.}{obsahuje len prirodzené čísla.}{obsahuje nekonečne veľa celých čísel.}{}{}}
\LANGes{
\Question{
Sea \( x\neq0 \). Completa la siguiente frase para obtener una proposición verdadera. El conjunto de soluciones de la inecuación \( \frac{|x|}x>2 \)}
{no contiene ningún número.}{contiene \( 2 \) números reales.}{contiene solo números naturales.}{contiene un número infinito de números enteros.}{}{}}
\End

\Data\ID{1003124205}
\Updated{2020-07-15 03:07:21}
\Level{B}
\SubArea{Absolute value}
\LANGen{
\Question{
Assuming \( x\in(4;7) \), the expression \( |x-4|-|x-7| \) can be written in the form:}
{\( 2x-11 \)}{\( -2x+11 \)}{\( 3 \)}{\( -11 \)}{}{}}
\LANGpl{
\Question{
Jeżeli \( x\in(4;7) \), to wyrażenie \( |x-4|-|x-7| \) można przedstawić w postaci:}
{\( 2x-11 \)}{\( -2x+11 \)}{\( 3 \)}{\( -11 \)}{}{}}
\LANGcs{
\Question{
Nechť \( x\in(4;7) \). Výraz \( |x-4|-|x-7| \) může být zapsán ve tvaru:}
{\( 2x-11 \)}{\( -2x+11 \)}{\( 3 \)}{\( -11 \)}{}{}}
\LANGsk{
\Question{
Nech \( x\in(4;7) \). Výraz \( |x-4|-|x-7| \) môže byť zapísaná v tvare:}
{\( 2x-11 \)}{\( -2x+11 \)}{\( 3 \)}{\( -11 \)}{}{}}
\LANGes{
\Question{
Suponiendo que \( x\in(4;7) \), la expresión \( |x-4|-|x-7| \) se puede escribir como:}
{\( 2x-11 \)}{\( -2x+11 \)}{\( 3 \)}{\( -11 \)}{}{}}
\End

\Data\ID{1003124207}
\Updated{2020-07-15 03:07:21}
\Level{B}
\SubArea{Absolute value}
\LANGen{
\Question{
On the real number line, the distance of a number \( x \) from the number \( -4 \) is given by:}
{\( |x+4| \)}{\( |x-4| \)}{\( |4x| \)}{\( |x|+4 \)}{}{}}
\LANGpl{
\Question{
Odległość liczby  \( x \) od liczby \( -4 \) na osi liczbowej jest równa:}
{\( |x+4| \)}{\( |x-4| \)}{\( |4x| \)}{\( |x|+4 \)}{}{}}
\LANGcs{
\Question{
Vzdálenost čísla \( x \) od čísla \( -4 \) dána na číselné ose se rovná:}
{\( |x+4| \)}{\( |x-4| \)}{\( |4x| \)}{\( |x|+4 \)}{}{}}
\LANGsk{
\Question{
Na číselnej osi, je vzdialenosť čísla \( x \) od čísla \( -4 \) daná:}
{\( |x+4| \)}{\( |x-4| \)}{\( |4x| \)}{\( |x|+4 \)}{}{}}
\LANGes{
\Question{
La distancia de un número \( x \) al número \( -4 \) en la recta numérica equivale a:}
{\( |x+4| \)}{\( |x-4| \)}{\( |4x| \)}{\( |x|+4 \)}{}{}}
\End

\Data\ID{1003124208}
\Updated{2020-07-15 03:07:21}
\Level{B}
\SubArea{Absolute value}
\LANGen{
\Question{
Assuming  \( -6 < x < 0 \), the   expression \( \frac{|x+6|-x+6}x \) is equal to:}
{\( \frac{12}x \)}{\( -\frac{12}x \)}{\( 2 \)}{\( 0 \)}{}{}}
\LANGpl{
\Question{
Dla każdej liczby x, spełniającej warunek  \( -6 < x < 0 \), wyrażenie \( \frac{|x+6|-x+6}x \) jest równe:}
{\( \frac{12}x \)}{\( -\frac{12}x \)}{\( 2 \)}{\( 0 \)}{}{}}
\LANGcs{
\Question{
Předpokládejme, že  \( -6 < x < 0 \). Výraz \( \frac{|x+6|-x+6}x \) se rovná:}
{\( \frac{12}x \)}{\( -\frac{12}x \)}{\( 2 \)}{\( 0 \)}{}{}}
\LANGsk{
\Question{
Predpokladajme, že  \( -6 < x < 0 \). Výraz \( \frac{|x+6|-x+6}x \) sa rovná:}
{\( \frac{12}x \)}{\( -\frac{12}x \)}{\( 2 \)}{\( 0 \)}{}{}}
\LANGes{
\Question{
Suponiendo \( -6 < x < 0 \), la expresión \( \frac{|x+6|-x+6}x \) equivale a:}
{\( \frac{12}x \)}{\( -\frac{12}x \)}{\( 2 \)}{\( 0 \)}{}{}}
\End

\Data\ID{1003124209}
\Updated{2020-07-15 03:07:21}
\Level{B}
\SubArea{Absolute value}
\LANGen{
\Question{
Which of the given inequalities holds for \( x=2\pi \)?}
{\( |x+1| > 5 \)}{\( |x-1| < 2 \)}{\( |x+3| \leq 4 \)}{\( |x-5| \geq 3 \)}{}{}}
\LANGpl{
\Question{
Wskaż nierówność, którą spełnia liczba  \( x=2\pi \):}
{\( |x+1| > 5 \)}{\( |x-1| < 2 \)}{\( |x+3| \leq 4 \)}{\( |x-5| \geq 3 \)}{}{}}
\LANGcs{
\Question{
Která z daných nerovnic platí pro \( x=2\pi \)?}
{\( |x+1| > 5 \)}{\( |x-1| < 2 \)}{\( |x+3| \leq 4 \)}{\( |x-5| \geq 3 \)}{}{}}
\LANGsk{
\Question{
Ktorá z daných nerovníc platí pre \( x=2\pi \)?}
{\( |x+1| > 5 \)}{\( |x-1| < 2 \)}{\( |x+3| \leq 4 \)}{\( |x-5| \geq 3 \)}{}{}}
\LANGes{
\Question{
¿Cuál de las desigualdades dadas es válida para  \( x=2\pi \)?}
{\( |x+1| > 5 \)}{\( |x-1| < 2 \)}{\( |x+3| \leq 4 \)}{\( |x-5| \geq 3 \)}{}{}}
\End

\Data\ID{1003124210}
\Updated{2020-07-15 03:07:21}
\Level{B}
\SubArea{Absolute value}
\LANGen{
\Question{
Which two numbers do both satisfy the equation \( |3x-3|=9 \)?}
{\( -2\text{, } 4 \)}{\( -4\text{, } 2 \)}{\( -5\text{, } 7 \)}{\( -7\text{, } 5 \)}{}{}}
\LANGpl{
\Question{
Liczbami spełniającymi równanie  \( |3x-3|=9 \) są:}
{\( -2\text{, } 4 \)}{\( -4\text{, } 2 \)}{\( -5\text{, } 7 \)}{\( -7\text{, } 5 \)}{}{}}
\LANGcs{
\Question{
Která dvě čísla splňují rovnici \( |3x-3|=9 \)?}
{\( -2\text{, } 4 \)}{\( -4\text{, } 2 \)}{\( -5\text{, } 7 \)}{\( -7\text{, } 5 \)}{}{}}
\LANGsk{
\Question{
Ktoré dve čísla spĺňajú rovnici \( |3x-3|=9 \)?}
{\( -2\text{, } 4 \)}{\( -4\text{, } 2 \)}{\( -5\text{, } 7 \)}{\( -7\text{, } 5 \)}{}{}}
\LANGes{
\Question{
¿Qué dos números satisfacen  la ecuación  \( |3x-3|=9 \)?}
{\( -2\text{, } 4 \)}{\( -4\text{, } 2 \)}{\( -5\text{, } 7 \)}{\( -7\text{, } 5 \)}{}{}}
\End

\Data\ID{1003187001}
\Updated{2020-07-15 03:07:21}
\Level{B}
\SubArea{Absolute value}
\LANGen{
\Question{
Let \( x\in(-\infty;-4] \). The value of the expression \( \left| |x|-4\right|-2|x-4|+|10-x| \) is equal to:}
{\( -2 \)}{\( -6 \)}{\( 2 \)}{\( 6 \)}{}{}}
\LANGpl{
\Question{
Dla \( x\in(-\infty;-4\rangle \). Wartość wyrażenia  \( \left| |x|-4\right|-2|x-4|+|10-x| \) jest równa:}
{\( -2 \)}{\( -6 \)}{\( 2 \)}{\( 6 \)}{}{}}
\LANGcs{
\Question{
Nechť \( x\in(-\infty;-4\rangle \). Hodnota výrazu \( \left| |x|-4\right|-2|x-4|+|10-x| \) se rovná:}
{\( -2 \)}{\( -6 \)}{\( 2 \)}{\( 6 \)}{}{}}
\LANGsk{
\Question{
Nech \( x\in(-\infty;-4\rangle \). Hodnota výrazu \( \left| |x|-4\right|-2|x-4|+|10-x| \) sa rovná:}
{\( -2 \)}{\( -6 \)}{\( 2 \)}{\( 6 \)}{}{}}
\LANGes{
\Question{
Sea \( x\in(-\infty;-4] \). El valor de la expresión  \( \left| |x|-4\right|-2|x-4|+|10-x| \) equivale a:}
{\( -2 \)}{\( -6 \)}{\( 2 \)}{\( 6 \)}{}{}}
\End

\Data\ID{1003187003}
\Updated{2020-07-15 03:07:21}
\Level{B}
\SubArea{Absolute value}
\LANGen{
\Question{
Simplifying the expression \( |3x-9|-|9-3x|+|-3x|-|-9| \) for \( x\in[3;9] \) you get:}
{\( 3x-9 \)}{\( -3x+9 \)}{\( 9x-27 \)}{\( 3x+9 \)}{}{}}
\LANGpl{
\Question{
Wyrażenie \( |3x-9|-|9-3x|+|-3x|-|-9| \) dla \( x\in\langle3;9\rangle \) jest równe:}
{\( 3x-9 \)}{\( -3x+9 \)}{\( 9x-27 \)}{\( 3x+9 \)}{}{}}
\LANGcs{
\Question{
Zjednodušením výrazu \( |3x-9|-|9-3x|+|-3x|-|-9| \) pro \( x\in\langle3;9\rangle \) obdržíme:}
{\( 3x-9 \)}{\( -3x+9 \)}{\( 9x-27 \)}{\( 3x+9 \)}{}{}}
\LANGsk{
\Question{
Zjednodušením výrazu \( |3x-9|-|9-3x|+|-3x|-|-9| \) pre \( x\in\langle3;9\rangle \) dostaneš:}
{\( 3x-9 \)}{\( -3x+9 \)}{\( 9x-27 \)}{\( 3x+9 \)}{}{}}
\LANGes{
\Question{
Simplificando la expresión \( |3x-9|-|9-3x|+|-3x|-|-9| \) para \( x\in[3;9] \) se obtiene:}
{\( 3x-9 \)}{\( -3x+9 \)}{\( 9x-27 \)}{\( 3x+9 \)}{}{}}
\End

\Data\ID{1003187004}
\Updated{2020-07-15 03:07:21}
\Level{B}
\SubArea{Absolute value}
\LANGen{
\Question{
Let \( x\in(-\infty;0) \). The expression \( \left|x-|x|\right| +\left|x+|x|\right|+1-x|x| \) is equal to:}
{\( (x-1)^2 \)}{\( (x+1)^2 \)}{\( x^2+1 \)}{\( 1-x^2 \)}{}{}}
\LANGpl{
\Question{
Dla  \( x\in(-\infty;0) \). Wartość wyrażenia \( \left|x-|x|\right| +\left|x+|x|\right|+1-x|x| \) jest równa:}
{\( (x-1)^2 \)}{\( (x+1)^2 \)}{\( x^2+1 \)}{\( 1-x^2 \)}{}{}}
\LANGcs{
\Question{
Nechť \( x\in(-\infty;0) \). Výraz \( \left|x-|x|\right| +\left|x+|x|\right|+1-x|x| \) se rovná:}
{\( (x-1)^2 \)}{\( (x+1)^2 \)}{\( x^2+1 \)}{\( 1-x^2 \)}{}{}}
\LANGsk{
\Question{
Nech \( x\in(-\infty;0) \). Výraz \( \left|x-|x|\right| +\left|x+|x|\right|+1-x|x| \) sa rovná:}
{\( (x-1)^2 \)}{\( (x+1)^2 \)}{\( x^2+1 \)}{\( 1-x^2 \)}{}{}}
\LANGes{
\Question{
Sea\( x\in(-\infty;0) \). La expresión  \( \left|x-|x|\right| +\left|x+|x|\right|+1-x|x| \) equivale a:}
{\( (x-1)^2 \)}{\( (x+1)^2 \)}{\( x^2+1 \)}{\( 1-x^2 \)}{}{}}
\End

\Data\ID{1003187104}
\Updated{2020-07-15 03:07:21}
\Level{B}
\SubArea{Absolute value}
\LANGen{
\Question{
Let \( a \) be a positive real number. Then \( |x| \leq a \)}
{if and only if \( -a \leq x \leq a \).}{if and only if \( x \leq a \).}{if and only if \( x \geq -a \).}{if and only if \( x < 0 \).}{}{}}
\LANGpl{
\Question{
Niech \( a \) będzie dowolną dodatnią liczbą rzeczywistą.  Dokończ, tak aby otrzymać zdanie prawdziwe: \( |x| \leq a \)}
{wtedy i tylko wtedy, gdy \( -a \leq x \leq a \).}{wtedy i tylko wtedy, gdy \( x \leq a \).}{wtedy i tylko wtedy, gdy \( x \geq -a \).}{wtedy i tylko wtedy, gdy \( x < 0 \).}{}{}}
\LANGcs{
\Question{
Nechť \( a \) je kladné reálné číslo. Pak \( |x| \leq a \)}
{právě tehdy, když \( -a \leq x \leq a \).}{právě tehdy, když \( x \leq a \).}{právě tehdy, když \( x \geq -a \).}{právě tehdy, když \( x < 0 \).}{}{}}
\LANGsk{
\Question{
Nech \( a \) je kladné reálne číslo. Pak \( |x| \leq a \)}
{práve vtedy, keď \( -a \leq x \leq a \).}{práve vtedy, keď \( x \leq a \).}{práve vtedy, keď \( x \geq -a \).}{práve vtedy, keď \( x < 0 \).}{}{}}
\LANGes{
\Question{
Sea \( a \) un número real positivo. Entonces \( |x| \leq a \)}
{si y solo si  \( -a \leq x \leq a \).}{si y solo si  \( x \leq a \).}{si y solo si \( x \geq -a \).}{si y solo si  \( x < 0 \).}{}{}}
\End

\Data\ID{1003187105}
\Updated{2020-07-15 03:07:21}
\Level{B}
\SubArea{Absolute value}
\LANGen{
\Question{
Let \( a \) be a positive real number. Then \( |x| \geq a \)}
{if and only if \( x \geq a \) or \( x \leq -a \).}{if and only if \( x \geq a \).}{if and only if \( x \leq -a \).}{if and only if \( x > 0 \).}{}{}}
\LANGpl{
\Question{
Niech \( a \) będzie dowolną liczbą rzeczywistą dodatnią. Dokończ, tak aby otrzymać zdanie prawdziwe:  \( |x| \geq a \)}
{wtedy i tylko wtedy, gdy \( x \geq a \) i \( x \leq -a \).}{wtedy i tylko wtedy, gdy \( x \geq a \).}{wtedy i tylko wtedy, gdy \( x \leq -a \).}{wtedy i tylko wtedy, gdy \( x > 0 \).}{}{}}
\LANGcs{
\Question{
Nechť \( a \) je kladné reálné číslo. Pak \( |x| \geq a \)}
{právě tehdy, když \( x \geq a \) nebo \( x \leq -a \).}{právě tehdy, když \( x \geq a \).}{právě tehdy, když \( x \leq -a \).}{právě tehdy, když \( x > 0 \).}{}{}}
\LANGsk{
\Question{
Nech \( a \) je kladné reálne číslo. Pak \( |x| \geq a \)}
{práve vtedy, keď \( x \geq a \) or \( x \leq -a \).}{práve vtedy, keď \( x \geq a \).}{práve vtedy, keď \( x \leq -a \).}{práve vtedy, keď \( x > 0 \).}{}{}}
\LANGes{
\Question{
Sea \( a \) un número real positivo. Entonces \( |x| \geq a \)}
{si y solo si \( x \geq a \) or \( x \leq -a \).}{si y solo si \( x \geq a \).}{si y solo si \( x \leq -a \).}{si y solo si \( x > 0 \).}{}{}}
\End

\Data\ID{1003187304}
\Updated{2020-07-15 03:07:21}
\Level{B}
\SubArea{Absolute value}
\LANGen{
\Question{
How many solutions does the equation \( \left| |x-4|-2\right|+2=0 \) have?}
{\( 0 \)}{\( 2 \)}{\( 4 \)}{\( 6 \)}{}{}}
\LANGpl{
\Question{
Ile rozwiązań ma równanie \( \left| |x-4|-2\right|+2=0 \)?}
{\( 0 \)}{\( 2 \)}{\( 4 \)}{\( 6 \)}{}{}}
\LANGcs{
\Question{
Kolik řešení má rovnice \( \left| |x-4|-2\right|+2=0 \)?}
{\( 0 \)}{\( 2 \)}{\( 4 \)}{\( 6 \)}{}{}}
\LANGsk{
\Question{
Koľko riešení má rovnica \( \left| |x-4|-2\right|+2=0 \)?}
{\( 0 \)}{\( 2 \)}{\( 4 \)}{\( 6 \)}{}{}}
\LANGes{
\Question{
¿Cuántas soluciones tiene la ecuación  \( \left| |x-4|-2\right|+2=0 \)?}
{\( 0 \)}{\( 2 \)}{\( 4 \)}{\( 6 \)}{}{}}
\End

\Data\ID{1003187405}
\Updated{2020-07-15 03:07:21}
\Level{B}
\SubArea{Absolute value}
\LANGen{
\Question{
Choose the expression which takes only negative values.}
{\( -|2-4x|-1 \)}{\( |-2x-5|-2 \)}{\( -|2-4x|+2 \)}{\( |2-x|-2 \)}{}{}}
\LANGpl{
\Question{
Tylko  wartości ujemne przyjmuje wyrażenie:}
{\( -|2-4x|-1 \)}{\( |-2x-5|-2 \)}{\( -|2-4x|+2 \)}{\( |2-x|-2 \)}{}{}}
\LANGcs{
\Question{
Vyberte výraz nabývající pouze záporných hodnot.}
{\( -|2-4x|-1 \)}{\( |-2x-5|-2 \)}{\( -|2-4x|+2 \)}{\( |2-x|-2 \)}{}{}}
\LANGsk{
\Question{
Vyberte výraz, ktorý nadobúda iba záporné hodnoty.}
{\( -|2-4x|-1 \)}{\( |-2x-5|-2 \)}{\( -|2-4x|+2 \)}{\( |2-x|-2 \)}{}{}}
\LANGes{
\Question{
Elige la expresión que tiene solo valores negativos.}
{\( -|2-4x|-1 \)}{\( |-2x-5|-2 \)}{\( -|2-4x|+2 \)}{\( |2-x|-2 \)}{}{}}
\End

\Data\ID{9000078505}
\Updated{2021-01-05 08:01:34}
\Level{B}
\SubArea{Absolute value}
\LANGen{
\Question{
Assuming \(x\in (0;\infty )\),
simplify the following expression. 
\[ 
 3x -|2x|-|- x|
\]}
{\(0\)}{\(2x\)}{\(3x\)}{\(4x\)}{}{}}
\LANGpl{
\Question{
Dla \(x\in (0;\infty )\),
uprość wyrażenie. 
\[ 
 3x -|2x|-|- x|
\]}
{\(0\)}{\(2x\)}{\(3x\)}{\(4x\)}{}{}}
\LANGcs{
\Question{
Nechť \(x\in (0;\infty )\). Výraz \[3x -|2x|-|- x|\] 
se rovná:}
{\(0\)}{\(2x\)}{\(3x\)}{\(4x\)}{}{}}
\LANGsk{
\Question{
Nech \(x\in(0;\infty)\). Výraz 
\[ 
 3x -|2x|-|- x|
\]
sa rovná:}
{\(0\)}{\(2x\)}{\(3x\)}{\(4x\)}{}{}}
\LANGes{
\Question{
Suponiendo que \(x\in (0;\infty )\),
simplifica la siguiente expresión.
\[ 
 3x -|2x|-|- x|
\]}
{\(0\)}{\(2x\)}{\(3x\)}{\(4x\)}{}{}}
\End

\Data\ID{9000078506}
\Updated{2021-01-05 08:01:17}
\Level{B}
\SubArea{Absolute value}
\LANGen{
\Question{
Assuming \(x\in (-\infty ;0)\),
simplify the following expression. 
\[ 
 3x -|2x|-|- x|
\]}
{\(6x\)}{\(4x\)}{\(2x\)}{\(0\)}{}{}}
\LANGpl{
\Question{
Dla  \(x\in (-\infty ;0)\),
uprość wyrażenie. 
\[ 
 3x -|2x|-|- x|
\]}
{\(6x\)}{\(4x\)}{\(2x\)}{\(0\)}{}{}}
\LANGcs{
\Question{
Nechť \(x\in (-\infty ;0)\). Výraz \[3x -|2x|-|- x|\] se 
rovná:}
{\(6x\)}{\(4x\)}{\(2x\)}{\(0\)}{}{}}
\LANGsk{
\Question{
Pre \(x\in(-\infty;0)\) je výraz \[3x-|2x|-|-x|\] rovný:}
{\(6x\)}{\(4x\)}{\(2x\)}{\(0\)}{}{}}
\LANGes{
\Question{
Suponiendo que \(x\in (-\infty ;0)\),
simplifica la siguiente expresión.
\[ 
 3x -|2x|-|- x|
\]}
{\(6x\)}{\(4x\)}{\(2x\)}{\(0\)}{}{}}
\End

\Data\ID{9000078507}
\Updated{2020-07-15 03:07:37}
\Level{B}
\SubArea{Absolute value}
\LANGen{
\Question{
Assuming \(x\in \left (-\frac{1}
{2};6\right )\),
simplify the following expression. 
\[ 
 3 -|6 - x| + |2x + 1|
\]}
{\(3x - 2\)}{\(x - 2\)}{\(3x + 10\)}{\(x + 8\)}{}{}}
\LANGpl{
\Question{
Dla \(x\in \left (-\frac{1}
{2};6\right )\),
uprość wyrażenie.
\[ 
 3 -|6 - x| + |2x + 1|
\]}
{\(3x - 2\)}{\(x - 2\)}{\(3x + 10\)}{\(x + 8\)}{}{}}
\LANGcs{
\Question{
Nechť \(x\in \left (-\frac{1}
{2};6\right )\). Výraz \[3 -|6 - x| + |2x + 1|\] 
se rovná:}
{\(3x - 2\)}{\(x - 2\)}{\(3x + 10\)}{\(x + 8\)}{}{}}
\LANGsk{
\Question{
Pre \(x\in(-\frac12;6)\) je výraz \[3-|6-x|+|2x+1|\] rovný:}
{\(3x - 2\)}{\(x - 2\)}{\(3x + 10\)}{\(x + 8\)}{}{}}
\LANGes{
\Question{
Suponiendo que \(x\in \left (-\frac{1}
{2};6\right )\),
simplifica la siguiente expresión.
\[ 
 3 -|6 - x| + |2x + 1|
\]}
{\(3x - 2\)}{\(x - 2\)}{\(3x + 10\)}{\(x + 8\)}{}{}}
\End

\Data\ID{9000078508}
\Updated{2021-01-05 10:01:26}
\Level{B}
\SubArea{Absolute value}
\LANGen{
\Question{
Assuming \(x\in (1;\infty )\),
simplify the following expression. 
\[ 
 3x -|2x + 1| + |x - 1|
\]}
{\(2x - 2\)}{\(4x - 2\)}{\(2x + 2\)}{\(2x\)}{}{}}
\LANGpl{
\Question{
Dla \(x\in (1;\infty )\),
uprość wyrażenie.
\[ 
 3x -|2x + 1| + |x - 1|
\]}
{\(2x - 2\)}{\(4x - 2\)}{\(2x + 2\)}{\(2x\)}{}{}}
\LANGcs{
\Question{
Pro \(x\in (1;\infty )\) 
je výraz \(3x -|2x + 1| + |x - 1|\) 
roven:}
{\(2x - 2\)}{\(4x - 2\)}{\(2x + 2\)}{\(2x\)}{}{}}
\LANGsk{
\Question{
Pre \(x\in (1;\infty )\) 
je výraz \(3x -|2x + 1| + |x - 1|\) 
rovný:}
{\(2x - 2\)}{\(4x - 2\)}{\(2x + 2\)}{\(2x\)}{}{}}
\LANGes{
\Question{
Suponiendo que \(x\in (1;\infty )\),
simplifica la siguiente expresión.
\[ 
 3x -|2x + 1| + |x - 1|
\]}
{\(2x - 2\)}{\(4x - 2\)}{\(2x + 2\)}{\(2x\)}{}{}}
\End

\Data\ID{9000078510}
\Updated{2021-01-05 08:01:21}
\Level{B}
\SubArea{Absolute value}
\LANGen{
\Question{
Assuming \(x\in (6;11)\),
simplify the following expression. 
\[ 
 3|x - 11|- 2|6 - x|
\]}
{\(- 5x + 45\)}{\(5x - 45\)}{\(x - 45\)}{\(x - 21\)}{}{}}
\LANGpl{
\Question{
Dla \(x\in (6;11)\),
uprość wyrażenie.
\[ 
 3|x - 11|- 2|6 - x|
\]}
{\(- 5x + 45\)}{\(5x - 45\)}{\(x - 45\)}{\(x - 21\)}{}{}}
\LANGcs{
\Question{
Pro \(x\in (6;11)\) 
je výraz \(3|x - 11|- 2|6 - x|\) 
roven:}
{\(- 5x + 45\)}{\(5x - 45\)}{\(x - 45\)}{\(x - 21\)}{}{}}
\LANGsk{
\Question{
Pre \(x\in (6;11)\) 
je výraz \(3|x - 11|- 2|6 - x|\) 
rovný:}
{\(- 5x + 45\)}{\(5x - 45\)}{\(x - 45\)}{\(x - 21\)}{}{}}
\LANGes{
\Question{
Suponiendo que \(x\in (6;11)\),
simplifica la siguiente expresión.
\[ 
 3|x - 11|- 2|6 - x|
\]}
{\(- 5x + 45\)}{\(5x - 45\)}{\(x - 45\)}{\(x - 21\)}{}{}}
\End

\Data\ID{9000081408}
\Updated{2020-07-15 03:07:37}
\Level{B}
\SubArea{Absolute value}
\LANGen{
\Question{
For \(x\in \mathbb{R}^{-}\) consider
expressions \(|x|\),
\(|- x|\),
\(-|x|\) and
\(- x\).
Which of these attains only negative values?}
{\(-|x|\)}{\(|x|\)}{\(|- x|\)}{\(- x\)}{}{}}
\LANGpl{
\Question{
Dla \(x\in \mathbb{R}^{-}\) zbadaj wyrażenie \(|x|\),
\(|- x|\),
\(-|x|\) i
\(- x\).
Które z nich ma tylko wartość ujemną?}
{\(-|x|\)}{\(|x|\)}{\(|- x|\)}{\(- x\)}{}{}}
\LANGcs{
\Question{
Jsou dány výrazy \(|x|\),
\(|- x|\),
\(-|x|\) a
\(- x\), kde
\(x\in \mathbb{R}^{-}\).
Vyberte variantu, v níž je uveden výraz nabývající pouze záporných
hodnot.}
{\(-|x|\)}{\(|x|\)}{\(|- x|\)}{\(- x\)}{}{}}
\LANGsk{
\Question{
Sú dané výrazy \(|x|\),
\(|- x|\),
\(-|x|\) a
\(- x\), kde
\(x\in \mathbb{R}^{-}\).
Vyberte variantu, v ktorom nižšie uvedený výraz nadobúda len záporné
hodnoty.}
{\(-|x|\)}{\(|x|\)}{\(|- x|\)}{\(- x\)}{}{}}
\LANGes{
\Question{
Para \(x\in \mathbb{R}^{-}\) considera las expresiones \(|x|\),
\(|- x|\),
\(-|x|\) y
\(- x\).
¿Cuál de ellas tendrá solo valores negativos?}
{\(-|x|\)}{\(|x|\)}{\(|- x|\)}{\(- x\)}{}{}}
\End

\Data\ID{1003049203}
\Updated{2020-07-15 03:07:12}
\Level{C}
\SubArea{Absolute value}
\LANGen{
\Question{
Identify which of the statements is false.}
{\( \forall a\text{, }b\in\mathbb{R}\colon |a+b|=|a|+|b| \)}{\( \forall a\text{, }b\in\mathbb{R}\colon |a\cdot b|=|a|\cdot|b| \)}{\( \forall a\in\mathbb{R}\text{, }b\in\mathbb{R}\setminus\{0\}\colon|\frac ab|=\frac{|a|}{|b|}  \)}{\( a\in\mathbb{R}\colon |a|=|-a| \)}{}{}}
\LANGpl{
\Question{
Określ, które z podanych  stwierdzeń jest fałszywe.}
{\( \forall a\text{, }b\in\mathbb{R}\colon |a+b|=|a|+|b| \)}{\( \forall a\text{, }b\in\mathbb{R}\colon |a\cdot b|=|a|\cdot|b| \)}{\( \forall a\in\mathbb{R}\text{, }b\in\mathbb{R}\setminus\{0\}\colon|\frac ab|=\frac{|a|}{|b|}  \)}{\( a\in\mathbb{R}\colon |a|=|-a| \)}{}{}}
\LANGcs{
\Question{
Vyberte nepravdivý výrok:}
{\( \forall a\text{, }b\in\mathbb{R}\colon |a+b|=|a|+|b| \)}{\( \forall a\text{, }b\in\mathbb{R}\colon |a\cdot b|=|a|\cdot|b| \)}{\( \forall a\in\mathbb{R}\text{, }b\in\mathbb{R}\setminus\{0\}\colon|\frac ab|=\frac{|a|}{|b|}  \)}{\( a\in\mathbb{R}\colon |a|=|-a| \)}{}{}}
\LANGsk{
\Question{
Ktoré tvrdenie nie je pravdivé?}
{\( \forall a\text{, }b\in\mathbb{R}\colon |a+b|=|a|+|b| \)}{\( \forall a\text{, }b\in\mathbb{R}\colon |a\cdot b|=|a|\cdot|b| \)}{\( \forall a\in\mathbb{R}\text{, }b\in\mathbb{R}\setminus\{0\}\colon|\frac ab|=\frac{|a|}{|b|}  \)}{\( a\in\mathbb{R}\colon |a|=|-a| \)}{}{}}
\LANGes{
\Question{
Identifica cuál de las proposiciones lógicas es falsa.}
{\( \forall a\text{, }b\in\mathbb{R}\colon |a+b|=|a|+|b| \)}{\( \forall a\text{, }b\in\mathbb{R}\colon |a\cdot b|=|a|\cdot|b| \)}{\( \forall a\in\mathbb{R}\text{, }b\in\mathbb{R}\setminus\{0\}\colon|\frac ab|=\frac{|a|}{|b|}  \)}{\( a\in\mathbb{R}\colon |a|=|-a| \)}{}{}}
\End

\Data\ID{1003187101}
\Updated{2020-07-15 03:07:21}
\Level{C}
\SubArea{Absolute value}
\LANGen{
\Question{
Which of the following relations is correct for all \( x \), \( y\in\mathbb{R} \)?}
{\( |x+y| \leq |x|+|y| \)}{\( |x+y|=|x|+|y| \)}{\( |x-y| < |x|-|y| \)}{\( |x-y|=|x|-|y| \)}{}{}}
\LANGpl{
\Question{
Która z poniższych zależności jest prawdziwa dla  \( x \), \( y\in\mathbb{R} \)?}
{\( |x+y| \leq |x|+|y| \)}{\( |x+y|=|x|+|y| \)}{\( |x-y| < |x|-|y| \)}{\( |x-y|=|x|-|y| \)}{}{}}
\LANGcs{
\Question{
Který z následujících vztahů platí pro všechna \( x \), \( y\in\mathbb{R} \)?}
{\( |x+y| \leq |x|+|y| \)}{\( |x+y|=|x|+|y| \)}{\( |x-y| < |x|-|y| \)}{\( |x-y|=|x|-|y| \)}{}{}}
\LANGsk{
\Question{
Ktorý z následujúcich vzťahov je správny pre všetky \( x \), \( y\in\mathbb{R} \)?}
{\( |x+y| \leq |x|+|y| \)}{\( |x+y|=|x|+|y| \)}{\( |x-y| < |x|-|y| \)}{\( |x-y|=|x|-|y| \)}{}{}}
\LANGes{
\Question{
¿Cuál de las siguientes relaciones es correcta para todo \( x \), \( y\in\mathbb{R} \)?}
{\( |x+y| \leq |x|+|y| \)}{\( |x+y|=|x|+|y| \)}{\( |x-y| < |x|-|y| \)}{\( |x-y|=|x|-|y| \)}{}{}}
\End

\Data\ID{1003187102}
\Updated{2020-07-15 03:07:21}
\Level{C}
\SubArea{Absolute value}
\LANGen{
\Question{
For \( x \), \( y\in\mathbb{R} \) consider \( |x+y|=|x|+|y| \).}
{The equality holds if and only if the sign of \( x \) and \( y \) is the same.}{The equality does not hold for any \( x \) and \( y \).}{The equality holds if and only if \( x \) and \( y \) are all positive.}{The equality holds if and only if \( x \) and \( y \) are both nonpositive.}{}{}}
\LANGpl{
\Question{
Dla \( x \), \( y\in\mathbb{R} \). Równość \( |x+y|=|x|+|y| \) jest prawdziwa:}
{Wtedy i tylko wtedy, gdy liczby \( x \) i \( y \) są tego samego znaku.}{Nie jest prawdziwa dla żadnego \( x \) i \( y \).}{Wtedy i tylko wtedy, gdy liczby \( x \) i \( y \) są  dodatnie.}{Wtedy i tylko wtedy, gdy liczby \( x \) i \( y \) są liczbami ujemnymi.}{}{}}
\LANGcs{
\Question{
Pro \( x \), \( y\in\mathbb{R} \) zvažte \( |x+y|=|x|+|y| \).}
{Rovnice platí právě tehdy, když čísla \( x \) a \( y \) mají stejné znaménko.}{Rovnice neplatí pro žádné \( x \) a \( y \).}{Rovnice platí právě tehdy, když \( x \) a \( y \) jsou všechna kladná.}{Rovnice platí právě tehdy, když \( x \) a \( y \) jsou obě nekladná.}{}{}}
\LANGsk{
\Question{
Pre \( x \), \( y\in\mathbb{R} \) zvážte \( |x+y|=|x|+|y| \).}
{Rovnica platí práve vtedy, keď znamienko \( x \) a \( y \) je rovnaké.}{Rovnice neplatí pre žiadne \( x \) a \( y \).}{Rovnice platí práve vtedy, keď \( x \) a \( y \) sú všetky kladné.}{Rovnica platí práve vtedy, keď \( x \) a \( y \) sú obidve nekladné.}{}{}}
\LANGes{
\Question{
Para \( x \), \( y\in\mathbb{R} \) considera \( |x+y|=|x|+|y| \).}
{La ecuación se cumple si y solo si el signo de \( x \) e \( y \) es el mismo.}{La ecuación no es válida para ningún \( x \) e \( y \).}{La ecuación se cumple si y solo si   \( x \) e \( y \) son ambos positivos.}{La ecuación se cumple si y solo si \( x \) e \( y \)  son ambos negativos.}{}{}}
\End

\Data\ID{1003187103}
\Updated{2020-07-15 03:07:21}
\Level{C}
\SubArea{Absolute value}
\LANGen{
\Question{
Find the relation that does not hold for any \( x \), \( y\in\mathbb{R} \).}
{\( \left| |x|-|y| \right| > |x+y| \)}{\( |xy|=|x| |y| \)}{\( \left|\frac xy \right|=\frac{|x|}{|y|}\text{, } y\neq0\text{ .} \)}{\( \left| (xy)^2 \right|=|xy|^2=(xy)^2 \)}{}{}}
\LANGpl{
\Question{
Wskaż relację, która nie zachodzi dla każdego \( x \), \( y\in\mathbb{R} \).}
{\( \left| |x|-|y| \right| > |x+y| \)}{\( |xy|=|x| |y| \)}{\( \left|\frac xy \right|=\frac{|x|}{|y|}\text{, } y\neq0\text{ .} \)}{\( \left| (xy)^2 \right|=|xy|^2=(xy)^2 \)}{}{}}
\LANGcs{
\Question{
Vyberte relaci, která neplatí pro žádné \( x \), \( y\in\mathbb{R} \).}
{\( \left| |x|-|y| \right| > |x+y| \)}{\( |xy|=|x| |y| \)}{\( \left|\frac xy \right|=\frac{|x|}{|y|}\text{, } y\neq0\text{ .} \)}{\( \left| (xy)^2 \right|=|xy|^2=(xy)^2 \)}{}{}}
\LANGsk{
\Question{
Vyberte reláciu, ktorá neplatí pre žiadne \( x \), \( y\in\mathbb{R} \).}
{\( \left| |x|-|y| \right| > |x+y| \)}{\( |xy|=|x| |y| \)}{\( \left|\frac xy \right|=\frac{|x|}{|y|}\text{, } y\neq0\text{ .} \)}{\( \left| (xy)^2 \right|=|xy|^2=(xy)^2 \)}{}{}}
\LANGes{
\Question{
Elige la relación que no se cumpla para ningún  \( x \), \( y\in\mathbb{R} \).}
{\( \left| |x|-|y| \right| > |x+y| \)}{\( |xy|=|x| |y| \)}{\( \left|\frac xy \right|=\frac{|x|}{|y|}\text{, } y\neq0\text{ .} \)}{\( \left| (xy)^2 \right|=|xy|^2=(xy)^2 \)}{}{}}
\End

\Data\ID{1003187106}
\Updated{2020-07-15 03:07:21}
\Level{C}
\SubArea{Absolute value}
\LANGen{
\Question{
Consider expressions \( |3x-12| \), \( 3|x|+12 \), \( |3x|-|-12| \), \( 3|x-4| \). Which of the expressions   has the greatest value if we substitute any \( x \) from the interval \( (0;+\infty) \)?}
{\( 3|x|+12 \)}{\( |3x-12| \)}{\( |3x|-|-12| \)}{\( 3|x-4| \)}{}{}}
\LANGpl{
\Question{
Które z podanych wyrażeń  \( |3x-12| \), \( 3|x|+12 \), \( |3x|-|-12| \), \( 3|x-4| \) ma największą wartość dla dowolnego  \( x \) z przedziału  \( (0;+\infty) \)?}
{\( 3|x|+12 \)}{\( |3x-12| \)}{\( |3x|-|-12| \)}{\( 3|x-4| \)}{}{}}
\LANGcs{
\Question{
Který z následujících výrazů \( |3x-12| \), \( 3|x|+12 \), \( |3x|-|-12| \), \( 3|x-4| \) nabývá největší hodnoty pro \( x\in(0;+\infty) \)?}
{\( 3|x|+12 \)}{\( |3x-12| \)}{\( |3x|-|-12| \)}{\( 3|x-4| \)}{}{}}
\LANGsk{
\Question{
Ktorý z nasledujúcich výrazov \( |3x-12| \), \( 3|x|+12 \), \( |3x|-|-12| \), \( 3|x-4| \) nadobúda najväčšie hodnoty pre \( x\in(0;+\infty) \)?}
{\( 3|x|+12 \)}{\( |3x-12| \)}{\( |3x|-|-12| \)}{\( 3|x-4| \)}{}{}}
\LANGes{
\Question{
Considera las expresiones  \( |3x-12| \), \( 3|x|+12 \), \( |3x|-|-12| \), \( 3|x-4| \). ¿Cuál de las expresiones tiene mayor valor si  \( x \) pertenece al intervalo \( (0;+\infty) \)?}
{\( 3|x|+12 \)}{\( |3x-12| \)}{\( |3x|-|-12| \)}{\( 3|x-4| \)}{}{}}
\End

\Data\ID{9000081406}
\Updated{2020-07-15 03:07:37}
\Level{C}
\SubArea{Absolute value}
\LANGen{
\Question{
For \(x\in \mathbb{R}\) find the correct
relationship between \(|x|\) 
and \(|- x|\).}
{\(|x| = |- x|\)}{\(|x| > |- x|\)}{\(|x| < |- x|\)}{None of them. The answer depends on the particular value of
\(x\).}{}{}}
\LANGpl{
\Question{
Dla \(x\in \mathbb{R}\) wskaż prawdziwą zależność pomiędzy \(|x|\) 
i \(|- x|\).}
{\(|x| = |- x|\)}{\(|x| > |- x|\)}{\(|x| < |- x|\)}{Żadna z nich. Odpowiedź zależy od poszczególnych wartości 
\(x\).}{}{}}
\LANGcs{
\Question{
Určete, jaký vztah platí mezi výrazy
\(|x|\) a
\(|- x|\), kde
\(x\in \mathbb{R}\).}
{\(|x| = |- x|\)}{\(|x| > |- x|\)}{\(|x| < |- x|\)}{Není možné jednoznačně určit. Záleží na hodnotě proměnné
\(x\).}{}{}}
\LANGsk{
\Question{
Určite, aký vzťah platí medzi výrazmi
\(|x|\) a
\(|- x|\), kde
\(x\in \mathbb{R}\).}
{\(|x| = |- x|\)}{\(|x| > |- x|\)}{\(|x| < |- x|\)}{Nie je možné jednoznačne určiť. Záleží na hodnote premennej
\(x\).}{}{}}
\LANGes{
\Question{
Para  \(x\in \mathbb{R}\) encuentra la relación correcta entre 
 \(|x|\) 
y \(|- x|\).}
{\(|x| = |- x|\)}{\(|x| > |- x|\)}{\(|x| < |- x|\)}{Ninguna de ellas. La respuesta depende del valor particular de \(x\).}{}{}}
\End

\Data\ID{9000081407}
\Updated{2020-07-15 03:07:37}
\Level{C}
\SubArea{Absolute value}
\LANGen{
\Question{
For \(x,y\in \mathbb{R}\) find the correct
relationship between \(|x - y|\) 
and \(|y - x|\).}
{\(|x - y| = |y - x|\)}{\(|x - y| > |y - x|\)}{\(|x - y| < |y - x|\)}{None of them. The answer depends on the particular values of
\(x\),
\(y\).}{}{}}
\LANGpl{
\Question{
Dla  \(x,y\in \mathbb{R}\) wskaż prawdziwą zależność pomiędzy \(|x - y|\) 
i \(|y - x|\).}
{\(|x - y| = |y - x|\)}{\(|x - y| > |y - x|\)}{\(|x - y| < |y - x|\)}{Żadna z nich. Odpowiedź zależy od poszczególnych wartości 
\(x\),
\(y\).}{}{}}
\LANGcs{
\Question{
Určete, jaký vztah platí mezi výrazy
\(|x - y|\) a
\(|y - x|\), kde
\(x,y\in \mathbb{R}\).}
{\(|x - y| = |y - x|\)}{\(|x - y| > |y - x|\)}{\(|x - y| < |y - x|\)}{Není možné jednoznačně určit. Záleží na hodnotě proměnných 
\(x\) a \(y\).}{}{}}
\LANGsk{
\Question{
Určte, aký vzťah platí medzi výrazmi
\(|x - y|\) a
\(|y - x|\), kde
\(x,y\in \mathbb{R}\).}
{\(|x - y| = |y - x|\)}{\(|x - y| > |y - x|\)}{\(|x - y| < |y - x|\)}{Nie je možné jednoznačne určiť. Záleží na hodnote premenných
\(x\) a \(y\).}{}{}}
\LANGes{
\Question{
Para \(x,y\in \mathbb{R}\) encuentra la relación correcta entre  \(|x - y|\) 
y \(|y - x|\).}
{\(|x - y| = |y - x|\)}{\(|x - y| > |y - x|\)}{\(|x - y| < |y - x|\)}{None of them. The answer depends on the particular values of
\(x\),
\(y\).}{}{}}
\End

\Data\ID{9000081409}
\Updated{2020-10-17 04:10:11}
\Level{C}
\SubArea{Absolute value}
\LANGen{
\Question{
Among expressions \(1 + |x|\),
\(|1 + x|\),
\(1 -|x|\) and
\(|1 - x|\) on the
set \(x\in (-\infty ;-1)\) 
find the expression which has smaller values than the other expressions from this list.}
{\(1 -|x|\)}{\(1 + |x|\)}{\(|1 + x|\)}{\(|1 - x|\)}{None of them.}{}}
\LANGpl{
\Question{
Wśród wyrażeń \(1 + |x|\),
\(|1 + x|\),
\(1 -|x|\) i
\(|1 - x|\) określonych na zbiorze
 \(x\in (-\infty ;-1)\) 
znajdź wyrażenie, które ma najmniejszą wartość.}
{\(1 -|x|\)}{\(1 + |x|\)}{\(|1 + x|\)}{\(|1 - x|\)}{Żadne z nich.}{}}
\LANGcs{
\Question{
Jsou dány výrazy \(1 + |x|\),
\(|1 + x|\),
\(1 -|x|\) a
\(|1 - x|\), kde
\(x\in (-\infty ;-1)\).
Který z uvedených výrazů má v daném definičním oboru nejmenší hodnotu.}
{\(1 -|x|\)}{\(1 + |x|\)}{\(|1 + x|\)}{\(|1 - x|\)}{Stejnou nejmenší hodnotu má více uvedených výrazů.}{}}
\LANGsk{
\Question{
Sú dané výrazy \(1 + |x|\),
\(|1 + x|\),
\(1 -|x|\) a
\(|1 - x|\), kde
\(x\in (-\infty ;-1)\).
Vyberte variantu, ktorá obsahuje výraz, ktorý má v danom oboru premennej
najmenšiu hodnotu.}
{\(1 -|x|\)}{\(1 + |x|\)}{\(|1 + x|\)}{\(|1 - x|\)}{Rovnakú najmenšiu hodnotu má viac uvedených výrazov.}{}}
\LANGes{
\Question{
Entre las expresiones \(1 + |x|\),
\(|1 + x|\),
\(1 -|x|\) y
\(|1 - x|\) donde \(x\in (-\infty ;-1)\) 
encuentra la expresión que tenga el valor más pequeño de todas ellas.}
{\(1 -|x|\)}{\(1 + |x|\)}{\(|1 + x|\)}{\(|1 - x|\)}{Ninguno de ellos.}{}}
\End
